\section{The placement of the localizers and their derivators}

\begin{claim}{H.1}{Proved}{The inclusions and the two poles}
Every modern fundamental localizer contains $\Winfty$. Among those satisfying Loc (5), exactly $\Winfty$ satisfies Loc ($7D_{\Winfty}$). The localizers satisfying Loc (6 b) have least member $\WZ$, but do not form the whole upper interval above $\WZ$. There are strict inclusions
\[
 \Winfty\subsetneq\W_+\subsetneq\WZ\subsetneq\W_0.
\]
Neither $\W_+$ nor $\W_0$ satisfies Loc (6 b). For each finite $n\geq1$, $\W_n$ and $\WZ$ are incomparable. No nontotal localizer satisfying Loc (6 b) contains $\W_1$.
\end{claim}
\begin{proof}
Minimality is \claimref{B.1}; the fibration classification, with its connectedness qualification, is \claimref{E.5}. The minimal reverse-excision class is \claimref{D.1}. The plus-construction inclusions and its failure of reverse excision are \claimref{F.9}; the $\W_0$ counterexample is \claimref{B.13}. Strictness of $\WZ\subset\W_0$ follows, for example, from $B\mathbb Z\to e$, since its source is connected and has first homology $\mathbb Z$. For $n\geq1$, the map $S^{n+1}\to *$ belongs to $\W_n$ and not to $\WZ$, while the nontrivial acyclic-group map $BG\to e$ of \claimref{B.10} belongs to $\WZ$ and not to $\W_n$. The sphere map means any specified simplex-category model; all such comparisons are weak homotopy equivalences. The final assertion is \claimref{D.3}. Thus the phrase ``the localizers above $\WZ$ with reversible cofibrations'' must be read as a subclass subject to that extra axiom, not as an inclusion-only characterization.
\end{proof}



\begin{claim}{H.2}{Proved}{Total asphericity and ordinary final images}
For a small category $A$, define $T_\W(A)$ by requiring that
\[
 B/(f_!1_A)\longrightarrow e\quad\text{belong to }\W
\]
for every small $B$ and every functor $f:A\to B$. Here $f_!$ is ordinary left Kan extension of presheaves along $f^{\op}$, and $B/H$ is the presheaf category of elements. Then $T_\W(A)$ holds if and only if every target $B$ of an ordinary final functor $A\to B$ is $\W$-aspherical. Final means that every $b\backslash A$ is nonempty and connected; no homotopical finality is implicit.
\end{claim}
\begin{proof}
For an arbitrary $f$, put $P_f(b)=\pi_0(b\backslash A)$. Precomposition by $\beta:b'\to b$ defines the presheaf action $P_f(b)\to P_f(b')$. The pointwise Kan formula is
\[
 (f_!1_A)(b)=\colim_{(b\backslash A)^{\op}}\{*\}=P_f(b).
\]
There is a functor $q_f:A\to B/P_f$ given by
$q_f(a)=(fa,[(a,\id_{fa})])$; on an arrow it uses its image under $f$. For an object $(b,c)$ of $B/P_f$, its comma category under $q_f$ is precisely the connected component $c$ of $b\backslash A$. It is nonempty and connected. Thus $q_f$ is ordinary final and the target in every Kan-extension test is a final image of $A$. Conversely, if $f$ is final, $P_f=1_B$, canonically, and $B/P_f\cong B$. This proves both directions, including empty categories. In particular, totality implies plain asphericity by the identity test. This verifies the exact final-image reformulation used in the attached TOTASPH return, with its presheaf variance specified.
\end{proof}

\begin{claim}{H.3}{Proved}{The entire interval is invisible to total asphericity}
For every small category $A$,
\[
 T_{\Winfty}(A)\quad\Longleftrightarrow\quad T_{\WZ}(A).
\]
Consequently these conditions are equivalent to $T_\W(A)$ for every $\Winfty\subseteq\W\subseteq\WZ$. Their intrinsic criterion is that $A$ be nonempty and connected and that, for every covariant functor $H:A\to\mathbf{Set}$, every connected component of $\El_A(H)$ be integrally acyclic. In this criterion ``integrally acyclic'' may equivalently be replaced by ``weakly contractible''.
\end{claim}
\begin{proof}
Denote the condition on all element categories by $\mathcal P_{\mathbb Z}(A)$, without imposing nonemptiness on an empty element category. Suppose first that $T_{\WZ}(A)$. The identity test makes $BA$ nonempty, connected, and integrally acyclic. For a covariant $H$ with $\El_A(H)$ nonempty and connected, adjoin an object $x$ with $\Hom(x,a)=H(a)$, no arrows from $A$ to $x$, and $\End(x)=\{\id\}$. The resulting category is the collage $A_H$ of \claimref{B.9}. The inclusion $A\to A_H$ is ordinary final: over $a\in A$ its comma category is $a\backslash A$, with initial object, and over $x$ it is $\El_A(H)$. Therefore $BA_H$ is integrally acyclic by \claimref{H.2}. The mapping-cone description of \claimref{B.9} and its reduced homology sequence show that $B\El_A(H)$ is integrally acyclic. For an arbitrary $H$, each connected component gives a subfunctor $H_E\subseteq H$, since every arrow preserves components. Applying the same argument to $H_E$ proves $\mathcal P_{\mathbb Z}(A)$.

The property is inherited by each component $E$ of each $\El_A(H)$. In fact, given $K:E\to\mathbf{Set}$, define
\[
 L(a)=\coprod_{h\in H_E(a)}K(a,h).
\]
For $v:a\to a'$, send the summand $K(a,h)$ by $K(v)$ into the summand indexed by $H_E(v)h$. The correspondence $((a,h),k)\mapsto(a,(h,k))$ is an isomorphism $\El_E(K)\cong\El_A(L)$, on both objects and arrows. This proves the inheritance assertion with the exact comparison, not merely an identification of homotopy types.

Let such an $E$ be connected and nonempty. The singleton functor makes $BE$ acyclic. Fix $e_0\in E$ and let $\Pi E$ be the groupoid obtained by formally inverting its arrows. The edge-path presentation identifies $G=\Aut_{\Pi E}(e_0)$ with $\pi_1(BE,e_0)$. For a subgroup $K\leq G$, put
\[
 H_K(e)=\Hom_{\Pi E}(e_0,e)/K,
\]
where $K$ acts by precomposition and arrows of $E$ act by postcomposition. These actions are bijective. Thus $\El_E(H_K)\to E$ has unique arrow lifts with either prescribed source or prescribed target. Its nerve realizes to the covering associated with $K$: edge paths lift uniquely, its monodromy is the transitive coset action, and a loop lifts to a loop at the identity coset exactly when its class belongs to $K$. In particular the covering is connected and its fundamental group is $K$. These are also the assertions of \citet[Theorem 1.38]{RootHatcher2002}, with the subgroup selected by the displayed basepoint.

Take $K$ to be the cyclic subgroup generated by an arbitrary element of $G$. The inherited property gives
\[
 0=H_1(B\El_E(H_K);\mathbb Z)=K_{\mathrm{ab}}=K.
\]
The first-homology identification is \citet[Theorem 2A.1]{RootHatcher2002}. Every element of $G$ is therefore trivial. Since $BE$ is a simply connected acyclic CW space, Hurewicz gives weak contractibility as in \claimref{B.6}. This proves that $\mathcal P_{\mathbb Z}$ and the corresponding weak-contractibility condition are equivalent.

Conversely, assume $A$ is nonempty and connected and satisfies this last condition. Its singleton element category shows that $BA$ is weakly contractible. For an ordinary final $f:A\to B$, every $b\backslash A$ is the element category of the covariant functor $a\mapsto\Hom_B(b,fa)$; it is nonempty and connected and hence weakly contractible. The final-variance Theorem A implies that the specified map $Bf$ is a weak equivalence. Its variance follows either by applying Theorem A to opposites or by the relative-Theorem-A property for $\Winfty$; reversing nerves does not change their realizations. Therefore $BB$ is weakly contractible. Claim \claimref{H.2} proves $T_{\Winfty}(A)$. The reverse implication to $T_{\WZ}(A)$ follows from inclusion of the localizers. Monotonicity of every test in \claimref{H.2} then proves the assertion for the entire interval. This independently verifies TOTASPH C.5.
\end{proof}

\begin{claim}{H.4}{Proved}{Above the one-type localizer, totality is plain asphericity}
If $\W_1\subseteq\W$, then $T_\W(A)$ is equivalent to $A\to e\in\W$ for every small $A$, including the empty category. No accessibility hypothesis is required.
\end{claim}
\begin{proof}
Only sufficiency needs proof. Let $f:A\to B$ be ordinary final and suppose $A\to e\in\W$. If $A$ is empty, finality forces $B$ empty. Otherwise finality gives a bijection on components and a surjection on each fundamental group. Here is the edge-path verification of the latter statement. For each object $b$ of $B$, choose an arrow $\beta_b:b\to f(a_b)$. For an arrow $v:b\to c$, the objects $(a_b,\beta_b)$ and $(a_c,\beta_cv)$ in $b\backslash A$ are joined by a zigzag. Applying $f$ to that zigzag identifies, in $\Pi B$, the arrow $\beta_cv\beta_b^{-1}$ with the image of a path in $\Pi A$. Multiplying these identities around any based loop cancels the $\beta_b$ factors and lifts the loop. Connectedness of the same comma categories shows that choices of a component over $b$ agree, and nonemptiness gives a component over every $b$; this proves the component bijection as well.

Choose one object $a_i$ in each component of $A$. Let $\mathcal S\subset\Pi A$ and $\mathcal T\subset\Pi B$ be the full subgroupoids on $a_i$ and $f(a_i)$ respectively. Their inclusions are equivalences, and $A\to\Pi A$, $B\to\Pi B$ are $\W_1$-equivalences by their edge-path presentations. The restricted functor $\psi:\mathcal S\to\mathcal T$ is bijective on objects and surjective on each endomorphism group. In particular it is a $\W$-cofibration by \claimref{C.11}. The ordinary pushout of $\mathcal T\leftarrow\mathcal S\to e$ is $e$: a compatible functor to any category must send all objects to the selected object and all arrows to its identity, since every arrow lifts under $\psi$. Cobase change along $\psi$ therefore sends the $\W$-map $\mathcal S\to e$ to $\mathcal T\to e$. The specified groupoid comparisons give $B\to e\in\W$. Claim \claimref{H.2} proves totality. For the total localizer the empty case is included automatically; for every other localizer, $\varnothing\to e$ is excluded by the argument in \claimref{D.3}. This verifies TOTASPH D.19.
\end{proof}

\begin{claim}{H.5}{Proved}{The precise boundary for all sieves to be cofibrations}
For a modern fundamental localizer $\W$, the following are equivalent: $\W_1\subseteq\W$; a category of homotopy type $S^2$ is $\W$-aspherical; every object-injective functor belongs to $\Cof_\W$; and every full sieve belongs to $\Cof_\W$.
\end{claim}
\begin{proof}
The equivalence of the first two conditions is \citet[Proposition 9.2.17]{RootCisinski2006}. Their implication to all object injections is \claimref{C.11}; a full sieve is object-injective. We prove the remaining implication by a complete finite example. Let $Q$ have objects $a,b,z,x$ and nonidentity arrows
\[
 r:a\to z,\quad v:b\to z,\quad
 \lambda,\mu:z\rightrightarrows x,\quad
 \alpha:a\to x,\quad\beta:b\to x,
\]
with the only nonidentity composites
$\lambda r=\mu r=\alpha$ and $\lambda v=\mu v=\beta$.
Its full subcategory $U=\{a,b,z\}$ is a sieve with terminal object $z$. The two full subcategories on $\{a,z,x\}$ and $\{b,z,x\}$ have initial objects and intersect in the parallel-arrow category $z\rightrightarrows x$, whose nerve is a circle. Their nerves cover $NQ$, so $BQ\simeq S^2$. Equivalently, this is the opposite of two terminal cones on that circle; reversing the nerve gives the same realization.

Collapsing $U$ to $e$ forces $r,v$ to be identities and hence forces $\lambda=\mu$. The pushout is the arrow category $[1]$. If $U\hookrightarrow Q$ is a $\W$-cofibration, applying its pushout functor to $U\to e\in\Winfty$ gives $Q\to[1]\in\W$, so $BQ$ is $\W$-aspherical. The program \texttt{code/sieve\_pushouts.py} constructs the categorical pushout by the defining coend relations, checks all compositions, and verifies the counts $(4,10)$ for $Q$, $(3,5)$ for $U$, and $(2,3)$ for its pushout. Its normalized integral chain calculations give Betti numbers $(1,0,1)$ for $Q$, with every nonzero Smith factor equal to $1$. The cone decomposition supplies the homotopy-type assertion, beyond what the homology calculation alone proves.
\end{proof}

\begin{claim}{H.6}{Refuted}{The pushout-product substitution for the manuscript's cylinders}
Let $n\geq2$ and let
$\Phi_n=(0\to1\leftarrow2\to\cdots\leftarrow2n)$.
If $S^2$ is not $\W$-aspherical, then the pushout-product of
$\{0\}\hookrightarrow[1]$ and $\{0,2n\}\hookrightarrow\Phi_n$
does not belong to $\Cof_\W$. In particular this fails for $\Winfty$ and for every nontotal localizer satisfying Loc (6 b). Moreover the only class of full-sieve $\W$-cofibrations stable under taking this pushout-product is a class consisting of isomorphisms. Products with the endpoint inclusion, as used in \claimref{B.4}, have no such defect.
\end{claim}
\begin{proof}
Write $P=[1]\times\Phi_n$, with lower and upper objects $l_k=(0,k)$ and $u_k=(1,k)$. Let $A\subset P$ contain every $l_k$ and only the two upper endpoints $u_0,u_{2n}$. It is a full sieve. Its inclusion $\alpha:A\hookrightarrow P$ is precisely the pushout-product in the statement. The nerve of $A$ is a tree, the lower alternating interval with two attached endpoint leaves. The category $P$ is a product of two contractible categories. Consequently $\alpha\in\Winfty$.

Put $C=Q^{\op}$ for the four-object category of \claimref{H.5}, retaining labels so its arrows are $\lambda,\mu:x\rightrightarrows z$, $r:z\to a$, $v:z\to b$, and their equal composites. Define $p:A\to C$ by $p(l_0)=z$, $p(u_0)=b$, and $p(w)=a$ at every other object. The arrows $l_0\to l_1$ and $l_0\to u_0$ go to $r$ and $v$; all remaining nonidentity arrows go to identities. Form the ordinary pushout $Q_0=P\amalg_A C$. We identify its crossing arrows to establish its homotopy type without replacing the categorical pushout by a homotopy pushout.

Its objects are the four objects of $C$ and the $u_k$, $1\leq k\leq2n-1$. It retains $C$ and the full subcategory of these upper objects, has no arrows from an upper object to $C$, and has crossing sets
\[
 \Hom_{Q_0}(c,u_k)
 =\int^{w\in A}\Hom_C(c,pw)\times\Hom_P(w,u_k).
\]
Composition is induced by that of $C$ and $P$. For $k\geq2$, every predecessor of $u_k$ in $A$ maps to $a$, and its predecessor category is nonempty and connected. The crossing presheaf is therefore $C(-,a)$. For $k=1$ its predecessor category has the four objects $l_0,l_1,l_2,u_0$ and arrows $l_0\to l_1$, $l_2\to l_1$, $l_0\to u_0$; its image is the fork $\Psi=(z\to a,z\to b)\subset C$. This functor to $\Psi$ is final: its three comma categories are respectively that connected tree, the arrow $l_2\to l_1$, and the singleton $u_0$. The fork inclusion $\Psi\to C$ is final as well. At $x$, its comma category has two objects over $z$ and one over each of $a,b$, with four crossing arrows, hence is a connected square; at the other objects its comma category has an initial object. Thus the coend for $k=1$ is the singleton presheaf on $C$. The transition from $u_2$ to $u_1$ is the unique map $C(-,a)\to1$; all transitions between the other representable weights are identities.

The full subcategory $R_0=C\cup\{u_1\}\subset Q_0$ is therefore the terminal cone on $C$. Define $\rho:Q_0\to R_0$ to fix this subcategory and to send every $u_k$, $k\geq2$, to $a$. On a crossing arrow represented by $h:c\to a$, use $h$; on an upper arrow use the corresponding identity of $a$ or its unique arrow to $u_1$. The coend description proves that these assignments respect composition. If $j:R_0\hookrightarrow Q_0$, then $\rho j=\id$ and there is a natural transformation $j\rho\Rightarrow\id_{Q_0}$, with components at $u_k$, $k\geq2$, the arrows represented by $\id_a$, and identity components elsewhere. Naturality on crossing arrows is exactly composition with $\id_a$; on $u_2\to u_1$ it is the unique-map transition above. Thus $BQ_0$ is contractible.

We use one general consequence of enough intrinsic cofibrations, already supplied by \claimref{B.4}. Every $i\in\Cof_\W\cap\W$ is couniversally in $\W$. Indeed factor an arbitrary $f:A\to D$ as $A\xrightarrow{k}D'\xrightarrow{w}D$ with $k\in\Cof_\W$ and $w\in\W$. Cobase change of $i$ along $k$ is in $\W$, since $k$ preserves $\W$. Cobase change of $w$ along $i$ is in $\W$, since $i$ preserves $\W$. The commuting square with left edge $w$ and these two maps makes the cobase change along $f$ belong to $\W$ by two out of three. If $\alpha$ were a $\W$-cofibration, its weak equivalence would therefore give $C\to Q_0\in\W$. Since $BQ_0$ is contractible, this would make $BC\simeq S^2$ aspherical, a contradiction.

There is also a literal weak-equivalence input exhibiting failure for $\Winfty$. Replace $C$ by $D=C^{\triangleright}$ with terminal object $s$, and use the composite $A\to C\to D$, which is a weak equivalence between contractible categories. The pushout $Q=P\amalg_A D$ retracts by the same formulas onto the full subcategory on $C,s,u_1$. The two added objects are terminal cones on $C$, with no arrows between them. Thus $BQ\simeq\Sigma BC\simeq S^3$, whereas $BP$ is contractible. Cobase change along $\alpha$ has sent the weak equivalence $A\to D$ to the non-weak-equivalence $P\to Q$.

All objects, arrows, and compositions for arbitrary $n$ have been specified by a finite poset, a four-object category, and explicit coend relations. The program checks the complete cases $n=2,3$, including every possible composite representative, all associativity equations, the pushout maps, the retraction, and its natural transformation. For $n=2$, $Q_0$ has $7$ objects and $25$ arrows and $Q$ has $8$ objects and $30$ arrows; their Betti numbers are respectively $(1,0,0,0,0)$ and $(1,0,0,1,0)$. For $n=3$ the respective counts are $(9,35)$ and $(10,40)$, with the same homology patterns. The uniform argument above proves the assertion for every $n\geq2$; these checks are certificates of finite instances, not an induction inferred from a sample.

Finally let $i:U\hookrightarrow X$ be a full sieve. If there is an arrow from $u\in U$ to $x\notin U$, selecting that arrow and using the characteristic functor $X\to[1]$ exhibit $\{0\}\hookrightarrow[1]$ as a retract of $i$ in the arrow category. The characteristic functor sends $U$ to $0$ and its complement to $1$; the sieve condition excludes arrows in the forbidden direction. Pushout-product with the fixed endpoint map preserves this retract diagram. Intrinsic cofibrations are closed under arrow retracts because the induced pushout comparison is an arrow retract of the corresponding preserved $\W$-map, and $\W$ is retract closed by \claimref{B.1}. Thus $i$ cannot belong to a class as in the statement. If there are no crossing arrows and $i$ is a nonisomorphism, then $X=U\amalg V$ with $V$ nonempty; its pushout-product has, for $v\in V$, the crossing arrow $(v,0)\to(v,1)$. Applying the first case to this new sieve contradicts closure under one further pushout-product. Hence every member of the class is an isomorphism. Claim \claimref{D.3} supplies the assertion about nontotal reversible localizers. This verifies, in the precise range $n\geq2$, the pushout-product objections in TOTASPH F.23--F.24 and F.31.
\end{proof}

\begin{claim}{H.7}{Proved}{Stability of the unpointed derivator}
Whenever the unpointed prederivator
\[
 \mathrm{HOT}_{\W}(I)=\Cat^{I^{\op}}[(\W^{I^{\op}})^{-1}]
\]
has the derivator structure under discussion, it is stable in Grothendieck's sense if and only if $\W=\Fl(\Cat)$. In particular none of the homology localizers for nonzero theories gives a stable unpointed derivator, although its equivalences admit the stable detection of Part F.
\end{claim}
\begin{proof}
Stability includes pointedness: initial and terminal objects agree. The images of $\varnothing$ and $e$ remain initial and terminal in $\Cat[\W^{-1}]$. To check this directly, the natural transformations from the constant $\varnothing$ functor to the identity and from the identity to the constant $e$ functor descend through localization: all three functors preserve $\W$, and naturality for inverted arrows follows from naturality for the original arrows. At $\varnothing$ and $e$ their respective components are identities. Naturality then says that every map from $\varnothing$, or to $e$, equals the displayed component, proving uniqueness as well as existence. If the derivator is pointed, its canonical map $\varnothing\to e$ is therefore invertible. Strong saturation \claimref{B.1} puts this specified functor in $\W$. The product and two-out-of-three argument of \claimref{D.3} makes $\W$ total.

Conversely, for the total class every morphism of every diagram category is inverted. The constant terminal-category diagram is terminal before localization and gives a contraction after groupoid completion: each object has its unique natural transformation to that terminal diagram. Thus every displayed value is equivalent to the terminal category. The resulting terminal derivator has its unique object both initial and terminal, and every square is both homotopy cartesian and homotopy cocartesian. It is stable. This concerns the unpointed $\mathrm{HOT}_{\W}$ requested here; neither a stable spectral test nor passage to a pointed or stabilized construction changes this conclusion about that specific derivator.
\end{proof}

\begin{claim}{H.8}{Proved}{The stronger separation from enough intrinsic fibrations}
Every nontotal modern fundamental localizer satisfying Loc (6 b) satisfies Loc (5) and fails Loc (7 F). Equivalently, Loc (6 b) and Loc (7 F) together force the total localizer, without any accessibility assumption. In particular every nontotal reversible localizer lies between $\WZ$ and $\W_0$, with a strict upper inclusion.
\end{claim}
\begin{proof}
Let $f:X\to Y$ be a map inverted by such a localizer, interpreted in spaces through the natural simplex-category models. Its pointed homotopy cofiber $C_f$ is $\W$-aspherical by \claimref{B.5}. If $f$ misses a component of $Y$, the cone component and an omitted component are distinct. Select one point in each and collapse all components into these two labels. This exhibits $S^0$ as a homotopy retract of $C_f$, contradicting \claimref{D.3} and retract closure.

Suppose instead that $\pi_0f$ is onto but not injective. Then $C_f$ is connected, and the ordinary homology exact sequence gives a surjection
\[
 H_1(C_f;\mathbb Z)\twoheadrightarrow
 \ker\bigl(H_0(X;\mathbb Z)\to H_0(Y;\mathbb Z)\bigr).
\]
The kernel is a nonzero free abelian group: over each component of $Y$ it is generated freely by differences from one selected component of $X$ above it. This also describes the kernel for infinite component sets. It has a quotient $\mathbb Z$. The resulting surjection $\pi_1C_f\twoheadrightarrow\mathbb Z$ is represented by a pointed map $h:C_f\to K(\mathbb Z,1)=S^1$. Choose a loop $s:S^1\to C_f$ whose image is $1$. The composite $hs$ has degree one and is pointed homotopic to the identity. Thus $S^1$ is a homotopy retract of $C_f$, again excluded by \claimref{D.3}. Homotopy retractions can be tested in $\Cat[\Winfty^{-1}]$ and then in $\Cat[\W^{-1}]$; strong saturation \claimref{B.1} transfers the resulting invertibility of $S^i\to *$ back to membership in $\W$. We have proved that every $f\in\W$ is a $\pi_0$-bijection, which is Loc (5).

Every $\W$-aspherical space $A$ is therefore connected and nonempty. No homomorphism $\pi_1A\to\mathbb Z$ is nonzero: otherwise its image, rescaled to $\mathbb Z$, and a chosen loop provide the same $S^1$ homotopy retraction. The free-subgroup argument in \claimref{C.4} now applies verbatim to show that $Y=BF_2$ is null for every such $A$; the constant-map map $Y\to\Map(A,Y)$ is an equivalence.

This last assertion needs no localization at a proper class of objects. Define directly
\[
 V_Y=\{u:A\to B:\Map(Bu,Y)\text{ is an equivalence of spaces}\}.
\]
It is a modern fundamental localizer. Weak saturation and the terminal axiom follow from the mapping-space functor. For relative Theorem A, use the natural comma-diagram comparisons $\hocolim_c B(A/c)\simeq BA$ and $\hocolim_c B(B/c)\simeq BB$ in the proof of \claimref{F.2}. Mapping into $Y$ turns those specified homotopy colimits into homotopy limits, so an equivalence for each comma map implies the required equivalence for $u$. Thus $\W^{\mathrm{asph}}\subseteq V_Y$.

If Loc (7 F) held, the accessibility-free direction of \claimref{E.2} would give $\W=\W^{\mathrm{asph}}\subseteq V_Y$. But the endomorphism $B\alpha:BF_2\to BF_2$ of \claimref{F.8} belongs to $\WZ\subseteq\W$ and does not belong to $V_Y$. Indeed, if precomposition by $B\alpha$ were an equivalence of $\Map(Y,Y)$, surjectivity on components would supply $r:Y\to Y$ with $rB\alpha\simeq\id_Y$. Injectivity on components, applied to $(B\alpha r)B\alpha\simeq B\alpha=\id_YB\alpha$, would then give $B\alpha r\simeq\id_Y$. This would make $B\alpha$ a homotopy equivalence, contradicting its explicitly nonsurjective fundamental-group homomorphism. Hence (7 F) fails. The total class satisfies both axioms directly, and the strict upper inclusion follows from the failure of reverse excision for $\W_0$ in \claimref{B.13}.
\end{proof}
